% Unit 7 map -- one page.
\documentclass{apchem-worksheet}

\apcunit{7}
\apcunittitle{Equilibrium}
\apcdoctitle{Unit Map}
\apcsubtitle{CED Unit 7 \textbullet\ 7--9\% of the AP Exam \textbullet\ 8 blocks + float + exam}

\begin{document}
\apctitleblock

\begin{apcnote}[How this unit is graded --- and what makes it different]
Unit exam \textbf{Mon Jan 25}: 30 multiple choice (\textbf{no calculator})
$+$ 1 short and 1 long free-response, 76 minutes, 44 points. This is one of
the two largest unit exams in the course.

Kinetics asked \emph{how fast}. This unit asks \emph{how far} --- and the
answer is a single number, $K$, that does not care about the path or the
rate. The central habit is comparing $Q$ with $K$: it decides which way a
system moves, whether it has arrived, and what happens when you disturb it.

Solubility (7.11--7.12) is the same machinery applied to a dissolving
solid, which is why it sits at the end rather than in a unit of its own.
\end{apcnote}

\section*{\sffamily Block calendar}

\renewcommand{\arraystretch}{1.3}
\noindent\begin{tabular}{@{}p{2.1cm}p{4.9cm}p{1.6cm}p{3.9cm}p{1.9cm}@{}}
\toprule
\textbf{Date} & \textbf{Topics} & \textbf{CED} & \textbf{Read (PDF pp.)} & \textbf{Practice}\\
\midrule
Mon Jan 4 & Dynamic equilibrium; reversibility & 7.1, 7.2 & \S13.1 \textbullet\ 651--654 & WS 1\\
Wed Jan 6 & $Q$ and $K$; calculating $K$ & 7.3, 7.4 & \S13.2--13.4 \textbullet\ 654--662 & WS 1\\
Fri Jan 8 & Magnitude and properties of $K$ & 7.5, 7.6 & \S13.3--13.5 \textbullet\ 658--671 & WS 2\\
Mon Jan 11 & ICE tables; the small-$x$ approximation & 7.7 & \S13.6 \textbullet\ 671--676 & WS 3\\
Wed Jan 13 & When the approximation fails; $K_p$ and $K_c$ & 7.7 & \S13.6, \S13.5 \textbullet\ 662--676 & WS 3\\
Fri Jan 15 & Representations; Le Ch\^atelier & 7.8, 7.9 & \S13.1, \S13.7 \textbullet\ 676--683 & WS 4\\
Mon Jan 18 & $Q$ versus $K$ and disturbances & 7.10 & \S13.7 \textbullet\ 676--683 & WS 4\\
Wed Jan 20 & Solubility equilibria; common ion & 7.11, 7.12 & \S16.1 \textbullet\ 805--812 & WS 5 (FRQ)\\
Fri Jan 22 & Reteach / float & --- & review sheet & ---\\
\textbf{Mon Jan 25} & \textbf{UNIT 7 EXAM} & all & review sheet & ---\\
\bottomrule
\end{tabular}

{\sffamily\footnotesize Dates from the co-op calendar; the unit resumes
after Christmas break. Unit 8 begins Wed Jan 27.}

\vspace{0.5em}
\section*{\sffamily By the end of this unit you can\ldots}

\begin{itemize}[leftmargin=*,itemsep=0.12em]
  \item Explain equilibrium as equal forward and reverse \emph{rates}, not
        equal amounts. \ced{7.1}, \ced{7.2}
  \item Write $Q$ and $K$ expressions, omitting pure solids and liquids.
        \ced{7.3}
  \item Calculate $K$ from equilibrium concentrations or pressures.
        \ced{7.4}
  \item Interpret the magnitude of $K$ and predict which side is favored.
        \ced{7.5}
  \item Adjust $K$ when a reaction is reversed, scaled, or added to
        another. \ced{7.6}
  \item Solve for equilibrium concentrations with an ICE table, and test
        any approximation. \ced{7.7}
  \item Interpret particulate and graphical representations of an
        equilibrium. \ced{7.8}
  \item Predict the effect of a disturbance using Le Ch\^atelier's
        principle. \ced{7.9}
  \item Justify that prediction quantitatively by comparing $Q$ with $K$.
        \ced{7.10}
  \item Relate $K_{sp}$ to molar solubility, and explain the common-ion
        effect. \ced{7.11}, \ced{7.12}
\end{itemize}

\begin{apctrap}[Where students lose points in this unit]
Saying equilibrium means equal concentrations; including pure solids or
liquids in $K$; forgetting to raise a concentration to its coefficient;
using the small-$x$ approximation without checking it, or checking it and
ignoring the result; claiming a catalyst or an inert gas shifts an
equilibrium; saying a concentration change alters $K$ (only temperature
does); and comparing $K_{sp}$ values for salts with different ion counts.
\end{apctrap}

\end{document}
